\section{Related Work}
\textbf{Vision-Language-Action Models.}
Large-scale vision-language models (VLMs) have significantly advanced multimodal learning by integrating image understanding and language reasoning~\cite{liu2024visual, bai2023qwen}. Extending these capabilities, VLA models~\cite{zitkovich2023rt,li2023vision} incorporate an action modality, enabling end-to-end visuomotor control. These models typically adopt large VLM backbones~\cite{touvron2023llama} and fine-tune them on robot data~\cite{o2024open}, with approaches varying from discretizing actions as language-like tokens~\cite{kim24openvla, chi2023diffusion} to incorporating specialized diffusion policy heads~\cite{black2024pi_0}. Despite their effectiveness in tasks like object retrieval and assembly~\cite{huang2024rekep, zhao2023learning}, VLA models demand substantial computation, making real-time deployment challenging, particularly in resource-constrained environments.

\textbf{Acceleration for Vision-Language Models.}
Inference acceleration has been extensively explored in vision-language models (VLMs) through quantization~\cite{10.1145/3664647.3680838}, pruning~\cite{lin2024mope}, and token-level techniques such as FastV~\cite{chen2025image}, SparseVLM~\cite{zhang2024sparsevlm}, ToMe~\cite{bolya2022token}, PuMer~\cite{cao2023pumer}, and MADTP~\cite{cao2024madtp}. These intra-frame strategies reduce redundancy within a \textbf{single image} but disregard the temporal and spatial structure essential for robotic tasks under closed-loop control. In the VLA domain, efficiency has been addressed through architectural modifications (\textit{e.g.}, RoboMamba~\cite{liu2024robomamba}, TinyVLA~\cite{wen2024tinyvla}), quantization-aware training (QAIL~\cite{park2024quantization}), and dynamic depth control (DeeR-VLA~\cite{DeeR-VLA}). While effective, these methods require re-training and lack generalizability. Recent high-frequency frameworks such as $\pi_0$-FAST~\cite{pertsch2025fast}, HiRT~\cite{zhang2024hirt}, and OpenVLA-OFT~\cite{kim2025fine} achieve higher control frequency via action chunking or asynchronous decoding. However, they continue to suffer from the \textbf{language model decoding} bottleneck, which dominates inference time. VLA-Cache addresses this gap by introducing a \textbf{cross-frame} token reuse strategy that accelerates inference without modifying the model or requiring additional training. Furthermore, it complements existing high-frequency VLA architectures by directly accelerating the language decoder bottleneck, offering a lightweight, plug-and-play solution for real-time robotic inference.
